% Template for PLoS
% Version 3.5 March 2018
%
% % % % % % % % % % % % % % % % % % % % % %
%
% -- IMPORTANT NOTE
%
% This template contains comments intended
% to minimize problems and delays during our production
% process. Please follow the template instructions
% whenever possible.
%
% % % % % % % % % % % % % % % % % % % % % % %
%
% Once your paper is accepted for publication,
% PLEASE REMOVE ALL TRACKED CHANGES in this file
% and leave only the final text of your manuscript.
% PLOS recommends the use of latexdiff to track changes during review, as this will help to maintain a clean tex file.
% Visit https://www.ctan.org/pkg/latexdiff?lang=en for info or contact us at latex@plos.org.
%
%
% There are no restrictions on package use within the LaTeX files except that
% no packages listed in the template may be deleted.
%
% Please do not include colors or graphics in the text.
%
% The manuscript LaTeX source should be contained within a single file (do not use \input, \externaldocument, or similar commands).
%
% % % % % % % % % % % % % % % % % % % % % % %
%
% -- FIGURES AND TABLES
%
% Please include tables/figure captions directly after the paragraph where they are first cited in the text.
%
% DO NOT INCLUDE GRAPHICS IN YOUR MANUSCRIPT
% - Figures should be uploaded separately from your manuscript file.
% - Figures generated using LaTeX should be extracted and removed from the PDF before submission.
% - Figures containing multiple panels/subfigures must be combined into one image file before submission.
% For figure citations, please use "Fig" instead of "Figure".
% See http://journals.plos.org/plosone/s/figures for PLOS figure guidelines.
%
% Tables should be cell-based and may not contain:
% - spacing/line breaks within cells to alter layout or alignment
% - do not nest tabular environments (no tabular environments within tabular environments)
% - no graphics or colored text (cell background color/shading OK)
% See http://journals.plos.org/plosone/s/tables for table guidelines.
%
% For tables that exceed the width of the text column, use the adjustwidth environment as illustrated in the example table in text below.
%
% % % % % % % % % % % % % % % % % % % % % % % %
%
% -- EQUATIONS, MATH SYMBOLS, SUBSCRIPTS, AND SUPERSCRIPTS
%
% IMPORTANT
% Below are a few tips to help format your equations and other special characters according to our specifications. For more tips to help reduce the possibility of formatting errors during conversion, please see our LaTeX guidelines at http://journals.plos.org/plosone/s/latex
%
% For inline equations, please be sure to include all portions of an equation in the math environment.
%
% Do not include text that is not math in the math environment.
%
% Please add line breaks to long display equations when possible in order to fit size of the column.
%
% For inline equations, please do not include punctuation (commas, etc) within the math environment unless this is part of the equation.
%
% When adding superscript or subscripts outside of brackets/braces, please group using {}.
%
% Do not use \cal for caligraphic font.  Instead, use \mathcal{}
%
% % % % % % % % % % % % % % % % % % % % % % % %
%
% Please contact latex@plos.org with any questions.
%
% % % % % % % % % % % % % % % % % % % % % % % %

\documentclass[10pt,letterpaper]{article}
\usepackage[top=0.85in,left=2.75in,footskip=0.75in]{geometry}

% amsmath and amssymb packages, useful for mathematical formulas and symbols
\usepackage{amsmath,amssymb}

% Use adjustwidth environment to exceed column width (see example table in text)
\usepackage{changepage}

% Use Unicode characters when possible
\usepackage[utf8x]{inputenc}

% textcomp package and marvosym package for additional characters
\usepackage{textcomp,marvosym}

% cite package, to clean up citations in the main text. Do not remove.
% \usepackage{cite}

% Use nameref to cite supporting information files (see Supporting Information section for more info)
\usepackage{nameref,hyperref}

% line numbers
\usepackage[right]{lineno}

% ligatures disabled
\usepackage{microtype}
\DisableLigatures[f]{encoding = *, family = * }

% color can be used to apply background shading to table cells only
\usepackage[table]{xcolor}

% array package and thick rules for tables
\usepackage{array}

% create "+" rule type for thick vertical lines
\newcolumntype{+}{!{\vrule width 2pt}}

% create \thickcline for thick horizontal lines of variable length
\newlength\savedwidth
\newcommand\thickcline[1]{%
  \noalign{\global\savedwidth\arrayrulewidth\global\arrayrulewidth 2pt}%
  \cline{#1}%
  \noalign{\vskip\arrayrulewidth}%
  \noalign{\global\arrayrulewidth\savedwidth}%
}

% \thickhline command for thick horizontal lines that span the table
\newcommand\thickhline{\noalign{\global\savedwidth\arrayrulewidth\global\arrayrulewidth 2pt}%
\hline
\noalign{\global\arrayrulewidth\savedwidth}}


% Remove comment for double spacing
%\usepackage{setspace}
%\doublespacing

% Text layout
\raggedright
\setlength{\parindent}{0.5cm}
\textwidth 5.25in
\textheight 8.75in

% Bold the 'Figure #' in the caption and separate it from the title/caption with a period
% Captions will be left justified
\usepackage[aboveskip=1pt,labelfont=bf,labelsep=period,justification=raggedright,singlelinecheck=off]{caption}
\renewcommand{\figurename}{Fig}

% Use the PLoS provided BiBTeX style
% \bibliographystyle{plos2015}

% Remove brackets from numbering in List of References
\makeatletter
\renewcommand{\@biblabel}[1]{\quad#1.}
\makeatother



% Header and Footer with logo
\usepackage{lastpage,fancyhdr,graphicx}
\usepackage{epstopdf}
%\pagestyle{myheadings}
\pagestyle{fancy}
\fancyhf{}
%\setlength{\headheight}{27.023pt}
%\lhead{\includegraphics[width=2.0in]{PLOS-submission.eps}}
\rfoot{\thepage/\pageref{LastPage}}
\renewcommand{\headrulewidth}{0pt}
\renewcommand{\footrule}{\hrule height 2pt \vspace{2mm}}
\fancyheadoffset[L]{2.25in}
\fancyfootoffset[L]{2.25in}
\lfoot{\today}

%% Include all macros below

\newcommand{\lorem}{{\bf LOREM}}
\newcommand{\ipsum}{{\bf IPSUM}}

\usepackage{color}
\usepackage{fancyvrb}
\newcommand{\VerbBar}{|}
\newcommand{\VERB}{\Verb[commandchars=\\\{\}]}
\DefineVerbatimEnvironment{Highlighting}{Verbatim}{commandchars=\\\{\}}
% Add ',fontsize=\small' for more characters per line
\usepackage{framed}
\definecolor{shadecolor}{RGB}{248,248,248}
\newenvironment{Shaded}{\begin{snugshade}}{\end{snugshade}}
\newcommand{\AlertTok}[1]{\textcolor[rgb]{0.94,0.16,0.16}{#1}}
\newcommand{\AnnotationTok}[1]{\textcolor[rgb]{0.56,0.35,0.01}{\textbf{\textit{#1}}}}
\newcommand{\AttributeTok}[1]{\textcolor[rgb]{0.77,0.63,0.00}{#1}}
\newcommand{\BaseNTok}[1]{\textcolor[rgb]{0.00,0.00,0.81}{#1}}
\newcommand{\BuiltInTok}[1]{#1}
\newcommand{\CharTok}[1]{\textcolor[rgb]{0.31,0.60,0.02}{#1}}
\newcommand{\CommentTok}[1]{\textcolor[rgb]{0.56,0.35,0.01}{\textit{#1}}}
\newcommand{\CommentVarTok}[1]{\textcolor[rgb]{0.56,0.35,0.01}{\textbf{\textit{#1}}}}
\newcommand{\ConstantTok}[1]{\textcolor[rgb]{0.00,0.00,0.00}{#1}}
\newcommand{\ControlFlowTok}[1]{\textcolor[rgb]{0.13,0.29,0.53}{\textbf{#1}}}
\newcommand{\DataTypeTok}[1]{\textcolor[rgb]{0.13,0.29,0.53}{#1}}
\newcommand{\DecValTok}[1]{\textcolor[rgb]{0.00,0.00,0.81}{#1}}
\newcommand{\DocumentationTok}[1]{\textcolor[rgb]{0.56,0.35,0.01}{\textbf{\textit{#1}}}}
\newcommand{\ErrorTok}[1]{\textcolor[rgb]{0.64,0.00,0.00}{\textbf{#1}}}
\newcommand{\ExtensionTok}[1]{#1}
\newcommand{\FloatTok}[1]{\textcolor[rgb]{0.00,0.00,0.81}{#1}}
\newcommand{\FunctionTok}[1]{\textcolor[rgb]{0.00,0.00,0.00}{#1}}
\newcommand{\ImportTok}[1]{#1}
\newcommand{\InformationTok}[1]{\textcolor[rgb]{0.56,0.35,0.01}{\textbf{\textit{#1}}}}
\newcommand{\KeywordTok}[1]{\textcolor[rgb]{0.13,0.29,0.53}{\textbf{#1}}}
\newcommand{\NormalTok}[1]{#1}
\newcommand{\OperatorTok}[1]{\textcolor[rgb]{0.81,0.36,0.00}{\textbf{#1}}}
\newcommand{\OtherTok}[1]{\textcolor[rgb]{0.56,0.35,0.01}{#1}}
\newcommand{\PreprocessorTok}[1]{\textcolor[rgb]{0.56,0.35,0.01}{\textit{#1}}}
\newcommand{\RegionMarkerTok}[1]{#1}
\newcommand{\SpecialCharTok}[1]{\textcolor[rgb]{0.00,0.00,0.00}{#1}}
\newcommand{\SpecialStringTok}[1]{\textcolor[rgb]{0.31,0.60,0.02}{#1}}
\newcommand{\StringTok}[1]{\textcolor[rgb]{0.31,0.60,0.02}{#1}}
\newcommand{\VariableTok}[1]{\textcolor[rgb]{0.00,0.00,0.00}{#1}}
\newcommand{\VerbatimStringTok}[1]{\textcolor[rgb]{0.31,0.60,0.02}{#1}}
\newcommand{\WarningTok}[1]{\textcolor[rgb]{0.56,0.35,0.01}{\textbf{\textit{#1}}}}


\usepackage{booktabs}
\usepackage{longtable}
\usepackage{array}
\usepackage{multirow}
\usepackage{wrapfig}
\usepackage{float}
\usepackage{colortbl}
\usepackage{pdflscape}
\usepackage{tabu}
\usepackage{threeparttable}
\usepackage{threeparttablex}
\usepackage[normalem]{ulem}
\usepackage{makecell}
\usepackage{xcolor}
\usepackage{booktabs}
\usepackage{longtable}
\usepackage{array}
\usepackage{multirow}
\usepackage{wrapfig}
\usepackage{float}
\usepackage{colortbl}
\usepackage{pdflscape}
\usepackage{tabu}
\usepackage{threeparttable}
\usepackage{threeparttablex}
\usepackage[normalem]{ulem}
\usepackage{makecell}
\usepackage{xcolor}



\usepackage{forarray}
\usepackage{xstring}
\newcommand{\getIndex}[2]{
  \ForEach{,}{\IfEq{#1}{\thislevelitem}{\number\thislevelcount\ExitForEach}{}}{#2}
}

\setcounter{secnumdepth}{0}

\newcommand{\getAff}[1]{
  \getIndex{#1}{Stats,Math,ERHS}
}

\providecommand{\tightlist}{%
  \setlength{\itemsep}{0pt}\setlength{\parskip}{0pt}}

\begin{document}
\vspace*{0.2in}

% Title must be 250 characters or less.
\begin{flushleft}
{\Large
\textbf\newline{Ten simple rules for selecting an R package} % Please use "sentence case" for title and headings (capitalize only the first word in a title (or heading), the first word in a subtitle (or subheading), and any proper nouns).
}
\newline
% Insert author names, affiliations and corresponding author email (do not include titles, positions, or degrees).
\\
Caroline J. Wendt\textsuperscript{\getAff{Stats}, \getAff{Math}},
G. Brooke Anderson\textsuperscript{\getAff{ERHS}}\textsuperscript{*}\\
\bigskip
\textbf{\getAff{Stats}}Department of Statistics, Colorado State University, Fort Collins,
Colorado, United States of America\\
\textbf{\getAff{Math}}Department of Mathematics, Colorado State University, Fort Collins,
Colorado, United States of America\\
\textbf{\getAff{ERHS}}Department of Environmental \& Radiological Health Sciences, Colorado
State University, Fort Collins, Colorado, United States of America\\
\bigskip
* Corresponding author: Brooke.Anderson@colostate.edu\\
\end{flushleft}
% Please keep the abstract below 300 words
\section*{Abstract}
R is an increasingly preferred software environment for data analytics
and statistical computing among scientists and practitioners. Packages
markedly extend R's utility and ameliorate inefficient solutions. We
outline ten simple rules for finding relevant packages and determining
which package is best for your desired use.

% Please keep the Author Summary between 150 and 200 words
% Use first person. PLOS ONE authors please skip this step.
% Author Summary not valid for PLOS ONE submissions.
\section*{Author summary}
Write the author summary here. Do we want to include and author summary?

\linenumbers

% Use "Eq" instead of "Equation" for equation citations.
\emph{Text based on plos sample manuscript, see
\url{http://journals.plos.org/ploscompbiol/s/latex}}

\hypertarget{disclaimer}{%
\section{Disclaimer?}\label{disclaimer}}

Do we need to include a disclaimer in the margin like the one from
{[}1{]} that states: ``\textbf{Competing Interests}: The authors have no
affiliation with GitHub, nor with any other commercial entity mentioned
in this article. The views described here reflect their own views
without input from any third party organization.''

\begin{itemize}
\tightlist
\item
  RStudio
\item
  ROpenSci
\item
  GitHub
\end{itemize}

I am an editor at ROpenSci, so we could mention that in the disclaimer.
I do not have any financial interests from that position (it's
volunteer, as are many journal editing positions in academics).

\hypertarget{funding-acknowledgment}{%
\section{Funding acknowledgment}\label{funding-acknowledgment}}

{[}Funding acknowledgement---Add in grant number for R25 and
acknowledgment of Honors program if appropriate{]}

\hypertarget{introduction}{%
\section{Introduction}\label{introduction}}

Computational reproducibility is surfacing as a central axiom in
academia, as researchers identify the need for transparency {[}2,3{]}.
The Methods section of scholarly work is not sufficient to convey the
complex data analysis methods and powerful tools of our time---data and
code must be shared. Some variability in scientific outcomes can be
attributed to differences in workflow and, today, the scientific
community deems the absence of automation to be irresponsible {[}4{]}.
Many traditional approaches that produce such inconsistencies also tend
to be at odds with productivity and collaboration. While standards for
disseminating and communicating computational science are evolving with
digitization, responsible researchers are adopting best practices so
others can more easily interpret, validate, and extend their published
knowledge {[}5{]}. These researchers rely on accessible and robust tools
to ensure their computations are reproducible.

As the R ecosystem---in which the life of modern data analysis
thrives---rapidly evolves alongside the burgeoning R community, R is
exhibiting sustained growth when compared to similar languages,
particularly in academia, healthcare, and government {[}6{]}. The open
source R language has become a dominant quantitative programming
environment in academic data analysis, enabling researchers to share
workflows and re-execute scripts within and across subsets of the
scientific community {[}4{]}. R is increasingly popular in computational
biology and bioinformatics, two of many disciplines generating
extensive, heterogenous, and complex data wanting for heavy-duty data
analysis tools that (ideally) support reproducibility {[}7,8{]}.

R was developed by statisticians and is collaboratively maintained by an
international group of core contributors {[}9{]}. Unlike several popular
proprietary languages such as MATLAB or SAS, R is freely available, open
source, and easily extensible software; the user can access, alter,
extend, and share code in various applications. Accordingly, a vibrant
community of R users has emerged, many of whom develop R packages to
complement and extend the functionality of base R. Contributed packages
comprise the bulk of enhancements made to the R environment {[}10{]}. In
fact, much of what we know about statistical methods and algorithms is
wrapped up in R packages---written and documented in various ways by R
users. There are numerous analogies in computing between programming and
culinary arts: recipe structures, coding cookbooks, and the like. To
conceptualize packages, imagine you are the chef, base R is the kitchen,
and packages are the special gadgets which allow you to cook and bake
new recipes. ``R package'' \emph{technically} refers to a collection of
R functions in a certain structure. The few tools that come with your
kitchen (e.g., \texttt{stats}) are indeed R packages, whereas ``R
extension'' describes the tools you add to your kitchen. It's important
to recognize that ``package'' is not synonymous with ``extension'';
however, since we don't want to keep distinguishing a kitchen sink from
a mixer, we use ``package'' in the colloquial sense. In essence, R
packages are coding delectables that enable the user to perform
practical tasks and solve problems with interesting techniques.

Are there R packages for wrangling and cleaning data frames, designing
interactive applications for data visualization, or performing
dimensionality reduction? Yes! How do you \emph{find} an R package to
help you train regression and classification models, assess the beta
diversity of a population, or analyze gene expression microarray data?
This answer is not as simple; there are tens of thousands of R packages.
As a natural consequence of the open source nature of R, there is
considerable variation in the quality of R packages and nontrivial
differences among those that provide similar tools.

Packages are essential to venturing beyond base R and, thus, quickly
become an integral aspect of advancing your R skills. Those who have
used R packages may know that, although leveraging existing tools can be
advantageous, the initial challenge of finding a suitable package for a
given task can obstruct potential benefits. The advanced R user---having
developed an intuition for their workflow---may be relatively confident
when searching for and selecting packages. By contrast, new R users who
are unfamiliar with the structure and syntax of the language may be
hindered by the process of finding packages because they do not know
where to search, what to look for, or how to sift through options. An
obstacle that characterizes learning R at the outset is the struggle to
(1) find a package to accomplish a particular task or solve a problem of
interest and (2) choose the best package to perform that task. Even so,
some obscure and complicated recipes make it difficult for an
experienced chef to select the best tools.

In both coding and life, we endeavor to make choices that optimize
outcomes. Just as one may go about shopping for shoes, deciding which
graduate program to pursue, or conducting a literature review, there is
a science behind selection. We inform our decisions by assessing,
comparing, and filtering options based on indicators of quality such as
utility, association, and reputation. Likewise, choosing an R package
requires attending to similar details. We outline ten simple rules for
finding and selecting R packages, so that you will spend less time
searching for the right tools and more time coding delightful recipes.

\hypertarget{rule-1-consider-your-purpose}{%
\section{Rule 1: Consider your
purpose}\label{rule-1-consider-your-purpose}}

Usually, there are several ways to accomplish a task while programming,
albeit some more elegant and efficient than others. Before looking for a
package to use for a task, determine whether you need one. Consider your
purpose by first identifying your goal then defining the scope and tasks
to achieve it. For some tasks and workflows, coding your own recipe with
existing tools is practical, while others benefit from new tools.

If the scope of your task is simple or reasonable, given your knowledge
and skills, using an R package may not be appropriate. There can be
advantages of coding in base R to complete your task or solve your
problem. First, when you code from scratch, you know precisely what you
are running; thus, your script may be easier to decipher and maintain
over time. Conversely, packages require you to rely on shared code with
features or underlying processes of which you may not be aware. Second,
although base R is relatively stable and slow to change, many R packages
evolve rapidly; unbeknown changes made to packages can cause defects in
portions of your code affected by these changes.

Packages are favorable in the same sense as kitchen tools: when a task
has a broad or complex scope beyond what you can (or desire to) attempt
from scratch. While there are ways to cook up an algorithm using
\texttt{for} loops and conditionals in base R, a relevant package may
accomplish your goal with less code and fewer bugs. The more reasonable
it is for a given task to be abstracted away from its context, the more
likely someone has generalized its themes, developed efficient
algorithms, and intuitively organized them to share with other R users.
Extensive tasks justify sophisticated frameworks with several functions
that form a cohesive package or even a suite of related packages. Data
manipulation is one such common task that has been streamlined by
packages such as \texttt{dplyr} and \texttt{tidyr} {[}11,12{]}, both
part of the \texttt{tidyverse} suite (see Table 1). Nevertheless, don't
be discouraged if you have a task that seems too unusual for a package.
There are indeed packages for seemingly singular tasks, which you may
favor over coding from scratch. Surprisingly, you can use an R package
to access the Twitter API, send emails from R, and there is even a
package for converting English letters to numbers as on a telephone
keypad {[}13--15{]}.

If you decide you need a package, next ask yourself: Which
functionalities in base R are restrictive in context of your task? Which
new functionalities would expand what you can do? If you identify
limitations of your current toolbox before searching for new tools, you
will be primed to recognize what you do and don't need. List
domain-specific keywords that describe what you are trying to do and how
you could do it, so you narrow your search. Identify the type of inputs
you have and envision working with them; contemplate the desired outputs
and corresponding format. Suppose you are using Bioconductor packages in
analyses and have outputs you want to visualize; you must consider that
the inputs may be of a certain class---namely, S4 objects---which impose
restrictions when creating graphics {[}16{]}. When you look for a
package to visualize the data, you will want to choose one that
addresses such restrictions. If, however, you don't quite know (or have
a difficult time articulating) what you need, don't hesitate to search
the internet for examples that capture your purpose. Sometimes, for
instance, it's helpful to search Google Images for plots and figures to
identify the terminology that describes what you are trying to do, how
people approach similar tasks, and potential challenges that may arise.

\hypertarget{rule-2-find-and-collect-options}{%
\section{Rule 2: Find and collect
options}\label{rule-2-find-and-collect-options}}

Tips and leads for finding R packages exist in a variety of places
online, in print, and elsewhere. If you've searched for packages a lot,
you likely have a shortlist of tried and true starting points---a perk
of being a long-time R user. On the other hand, new R users who haven't
developed an internal compass to navigate the initial ins and outs of
finding a package, don't know where to start looking for a package to
suit their task. Here are a few directions to head in: You can discover
new packages any time you learn R-related topics, browse the internet,
or connect with the international community of R users.

\textbf{Learn}

When learning how to program in R, you are typically introduced to some
of the most common packages, which tend to have more general purposes
(\textbf{Table 1}). In addition, reputable online tutorials, courses,
and books are helpful resources for acquiring knowledge about packages
that are versatile and reliable---many of which are short, accessible,
and either affordable or offered at
\href{https://committedtotape.shinyapps.io/freeR/}{no cost} to the
learner. We recommend online R programming courses such as those through
\href{https://www.coursera.org/learn/r-programming}{Coursera} and
\href{https://www.codecademy.com/learn/learn-r}{Codeacademy} for
interactive learning and R book series including the
\href{https://rstudio.com/resources/books/}{RStudio books} and Springer
titles for further reading.

\textbf{Browse}

The solution-seeking tactics we employ for many tasks nowadays may lead
you to think that finding R packages relies heavily on internet search
queries. Indeed, search engines such as Google return ample pages
related to anything \texttt{"\ldots{}in\ R"}. Keep in mind, however,
this approach sometimes leads to frustration and confusion when
attempting to find a package tailored to your purpose (see Rule 1). For
a more directed approach, you can search the package lists on
repositories including the Comprehensive R Archive Network (CRAN),
Bioconductor, and rOpenSci or browse packages available on code-sharing
websites, with language filters for finding repositories that heavily
rely on R code, like GitHub and GitLab---all of which will be further
discussed in Rule 3.

That said, in many cases, you may find it more helpful to start from a
curated list of R packages on a particular topic. There are several
available. First, \href{https://cran.r-project.org/web/views/}{CRAN Task
Views} are concentrated topics from certain disciplines and
methodologies related to statistical computing that categorize R
packages by the tasks they perform (e.g., Econometrics, Genetics,
Optimization, Spatial). In the HTML version, you can browse alphabetized
subcategories within each Task View and read concise descriptions to
find tools with specific functions. Alternatively, you can access Task
Views directly from the R console with \texttt{ctv::CRAN.views()}, which
is compact notation of the form \texttt{package::function()}. To date,
there are 41 Task Views that collectively contain thousands of packages
which are curated and regularly tested. Moreover, CRAN Task Views
provide tools that enable you to automatically install all packages
within a targeted area of interest. Ultimately, by providing task-based
organization, easy simultaneous installation of related packages,
meta-information, ensured maintenance, and quality control, CRAN Task
Views address several major user-end issues that have arisen due to the
sheer quantity of available packages {[}17{]}.

The computational biology and bioinformatics analog of CRAN Task Views,
\href{https://www.bioconductor.org/packages/release/BiocViews.html\#___Software}{biocViews},
is a list of packages, available through the Bioconductor website, that
provide tools for analyzing biological data. Your keywords from Rule 1
will come in handy here; for instance, you can search ``RNA-seq'' to
find all packages related to high-throughput sequencing methods for
analyzing RNA populations. You can also browse the biocViews tree menu
to narrow your search based on which general and specific categories
describe your task. The list includes each package's name linked to
information and documentation, its maintainer, a short description, and
its rank relative to other Bioconductor packages.

Other curated collections and topically-linked lists of packages are
also available.
\href{http://dirk.eddelbuettel.com/cranberries/index.html}{CRANberries},
for example, is a hub of information about new, updated, and removed
packages from CRAN. Another place to find well-maintained tools is
\href{https://ropensci.org/packages/}{rOpenSci packages}. These
filterable and searchable peer-reviewed (see Rule 6) R packages are
organized by name, maintainer, description, and status markers based on
activity, association, and review.

\textbf{Embrace the community}

An inclusive and collaborative community is an overlooked, yet integral,
aspect of a software's success {[}18{]}. A defining feature of R is the
enthusiasm of its users and developers alike. The R community has a
widespread internet presence across various platforms; however, members
are markedly active on Twitter, a place where R users seek help, share
ideas, and stay informed on \texttt{\#rstats} happenings, including
releases of new packages {[}19{]}. R-related blogs serve as another
informal, up-to-date, and more detailed avenue for communicating and
promoting R-related information (\textbf{Table ?}). For example, Joseph
Rickert, Ambassador at Large for RStudio, writes monthly posts on the
\href{https://rviews.rstudio.com/}{R Views} blog highlighting
interesting new R packages. Rickert also features special articles about
recently released packages and lists of top packages within certain
categories, including Computational Methods, Data, Machine Learning,
Medicine, Science, Statistics, Time Series, Utilities, and
Visualization.

You can also attend---or access content from---conferences to learn
about improved, recently released, or up-and-coming R package
developments and applications. Two large annual R conferences are
\href{https://rstudio.com/conference/}{rstudio::conf} for industry and
\href{https://www.r-consortium.org/}{useR!} (not exclusively) for
academia. As supportive and innovative subsets of the R community, some
boutique industries that use R host smaller conferences:
\href{https://rinpharma.com/}{R/Pharma} (pharmaceutical development),
\href{https://www.rinfinance.com/}{R/Finance} (applied finance), and
\href{https://www.bioconductor.org/help/events/}{BioC} (open source
software for bioinformatics). Conferences in your field may foster
connections with fellow scientists who use R for similar tasks and help
you collect information about packages related to your expertise. Talks
and presentations at conferences are often recorded and made available
online for playback if you can't attend in person or wish to revisit
content from past conferences. For instance, R Studio has a
\href{https://rstudio.com/resources/rstudioconf-2020/}{webpage} devoted
to videos of past conferences by year and the talks from useR! are
posted on the R Consortium
\href{https://www.youtube.com/channel/UC_R5smHVXRYGhZYDJsnXTwg/featured}{YouTube
channel}. Of late, the R community and rOpenSci have popularized
nontraditional Unconferences and R Collaboratives at which people
interact and share ideas in a collegial and unstructured atmosphere; as
you may imagine, R packages are a hot topic of conversation at these
events.

Don't forget, a developer of an R package may intend for it to be
private (exclusively for personal or professional use) or public (at no
cost and available for use by anyone) {[}20{]}. If your task is specific
to a line of research, consult colleagues to see if they have a relevant
(private) package they would be willing to share.

\hypertarget{rule-3-check-how-its-shared}{%
\section{Rule 3: Check how it's
shared}\label{rule-3-check-how-its-shared}}

Packages can be shared through a variety of platforms. Repositories are
a primary way in which developers share packages for both public and
private use, but there are alternatives. As far as R packages are
concerned, a repository is essentially a warehouse for tools before they
are shipped to your kitchen; in computing terms, a repository is
analogous to a cloud because it is a central location in which data is
stored and managed. Some repositories impose vetting mechanisms that
tame unwieldy aspects of the R ecosystem by regularly checking
underlying code and managing corresponding webs of dependencies. Two
traditional repositories for R packages are CRAN and Bioconductor;
however, there are lesser-known remote repositories that have unique
properties. Developers can also share their R packages through large
code-sharing sites, including public version control platforms like
GitHub and GitLab; these have fewer restrictions on the format or
content of shared code compared to repositories. At the least public
level, in lieu of making packages accessible to everyone via internet
repositories, some developers share their code in zipped files directly
with collaborators (see Rule 2).

The CRAN package repository is the most established and primary source
from which you can install R packages. CRAN hosts a huge collection of R
packages (over 16,000 as of August 2020) on almost any conceivable
topic, from random forests to multidimensional maps {[}21,22{]}. Due to
its longevity and historical role, Rickert asserts that ``CRAN is the
greatest open source repository for statistical computing knowledge in
the world'' {[}20{]}. A key advantage of using a package from CRAN is
that the repository integrates easily with your base R installation
{[}23{]}. You can simply use the \texttt{install.packages()} function to
install a package from CRAN; source and/or binary code are automatically
saved to your computer in a designated package library {[}24{]}. When
you want to use a particular package, you load it to your R session via
the \texttt{library()} function from base R. The R Foundation manages
CRAN and imposes some strict regulatory practices for the selection and
maintenance of the packages they host, particularly as standards pertain
to the structure of a package. CRAN is more or less mechanical in their
approach to accepting and keeping a package as it evolves in the R
ecosystem. CRAN is not as concerned with content or substance because
they do not have a peer review process (see Rule 3 and Rule 6). At
worst, a CRAN package might have a solid framework with helpfiles for
each function, but its underlying code might not necessarily do what it
claims. A package must pass a series of initial stability tests in
accordance with the CRAN Repository Policy before obtaining publication
privileges {[}25{]}. A number of CRAN packages make impactful
contributions to certain subject areas including statistical computing
and graphics; however, some CRAN packages, such as
\texttt{weathermetrics} for converting between weather metrics, are more
convenient than substantial {[}26{]}. Volunteer CRAN maintainers
actively monitor portions of code (e.g., vignettes, helpfiles, test
directories) for source and contributed packages to ensure compatibility
with the latest version of R and dependencies; they modify or remove
packages that do not uphold these publication quality standards.

If you're looking for tools specifically designed for high-throughput
genomic data, the Bioconductor repository is a more specialized
repository than CRAN that's worth investigating. The Bioconductor
project was motivated by a need for transparent, reproducible, and
efficient software in computational biology and bioinformatics and
supports the integration of computational rigor and reproducibility in
research on biological processes {[}7{]}. Bioconductor packages
facilitate the analysis and comprehension of biological data and help
users solve problems that arise when working with high-throughput
genomic data such as those related to microarrays, sequencing, flow
cytometry, mass spectrometry, and image analysis {[}27{]}. Bioconductor
boasts a modularized design, wherein data structures, functions, and the
packages that contain them have distinct roles that are accompanied by
thorough documentation. Accessibility is a pillar of the Bioconductor
project; all forms of documentation, including courses, vignettes, and
interactive documents, are curated for individuals with expertise in
adjacent disciplines and minimal experience with R (see Rule 4) {[}7{]}.
Bioconductor packages integrate well with base R by different means than
CRAN packages; you should use \texttt{BiocManager::install()} to install
Bioconductor packages. Bioconductor has a structured release schedule
and thus a unique package management system which automatically loads
your Bioconductor packages in versions that are all from the same
release cycle and compatible with your current version of R. The R
environment and its package system is fundamental to the implementation
of Bioconductor's interoperable and object-oriented (S4) infrastructure.
The tasks for which Bioconductor packages provide tools generally
involve large and complex datasets that need to leverage R's object
oriented programming functionalities---a set of specific object classes
in particular (e.g., ExpressionSet class) {[}28{]}. This helps, among
other things, keep metadata properly aligned with sample data. Similar
to CRAN, Bioconductor has strict criteria for package submissions: a
package must be relevant to high-throughput genomic analysis,
interoperable with other Bioconductor packages, well-documented,
supported in the long-term, and comply with additional package
guidelines {[}29{]}. Unlike CRAN, however, to ensure Bioconductor
software maintains a centralized content focus, their packages are
coordinated and peer-reviewed.

Additionally, there are smaller, specialized repositories that share R
packages. The non-profit organization, rOpenSci, for example, runs a
repository as part of their commitment to promote open science practices
through technical and social infrastructure for the R community
{[}30{]}. The repository only includes packages that have either been
developed by staff to fill a gap in infrastructure or have passed their
open review process. These packages are within the scope and aims
defined by rOpenSci, including those that focus on data extraction,
workflow automation, and field or laboratory reproducibility tools
{[}31{]}. R package maintainers can even create and host their own
personalized repository using the \texttt{drat} (Drat R Archive
Template) package, which enables developers to design individual
repositories and suites of coordinated repositories for packages that
are stored in and/or distributed through platforms like GitHub {[}32{]}.
In some cases, this might be done to host a package that does not meet
certain restrictions from other repositories (e.g., to share data
through a package that is larger than CRAN's size limit for packages
{[}33{]}). There are other platforms strictly for the development,
rather than distribution, of R packages such as R-Forge and Omegahat,
which are beyond the scope of this paper {[}34,35{]}.

While there are evident benefits of exclusive repositories, there are
also considerable advantages to hosting a package on online hosting
platforms for git repositories like GitHub or GitLab {[}20,36{]}. These
platforms can be used for production of R packages, including
large-scale packages. Package developers can post the directory with all
source code for their package as they develop it, from the earliest
stages. The platforms allow for continuous integration and continuous
deployment for testing data and code to ensure a stable end-product. One
particularly popular platform when it comes to sharing R packages is
GitHub. An increasing number of packages are hosted on GitHub during the
development stages, and so it will often be the first place you can find
up-and-coming packages (see Rule 2). Further, even if developers choose
not to distribute their package through GitHub, the stable release
versions of such packages are often published on CRAN or Bioconductor
{[}37{]}, and for these GitHub can be a place to check out planned
changes for the next stable version that the developer plans to publish.
GitHub provides R users with open access to package code, a timeline of
help resources (see Rule 4), and a direct line of communication to
developers. It is simple to install the latest version of a package you
want to use from GitHub via \texttt{devtools::install\_github()};
however, the decentralized nature of GitHub is not conducive to a tool
that automatically locates and installs corresponding dependencies
{[}38{]}. In fact, the rapid uptick in package development and ensuing
inter-repository dependencies has sparked an ongoing debate on whether
regulated repositories such as CRAN and Bioconductor are preferable to
distribution platforms like GitHub for this reason {[}20,36,39{]}.
Further, when a package is installed from GitHub, typically only the
source code is available, not binaries pre-built for your computers
operating systems. This means that the installation process will need to
build the package from source on your computer, which may require
certain compilers and other tools (e.g., XCode for macOS, Rtools for
Windows). Both R users and package developers benefit from interactive
feedback channels through GitHub Issues and the Star rating system.

When you are searching for and selecting packages, it is important to
note that a single package is often available in multiple places.
Packages that are on CRAN, Bioconductor, or rOpenSci are usually on
GitHub. There is an unofficial read-only GitHub mirror at CRAN,
\href{https://github.com/cran}{METACRAN}, in which \emph{all} CRAN
packages have their own GitHub repositories with all versions, original
dates, and authors. Further, package developers often use their personal
GitHub accounts to publicize a package while it is being developed and
continue to host development versions after it has been published on
formal repositories such as CRAN. In this case, CRAN would host a
``stable'' version of the package, which may simultaneously be in the
``master'' branch of the developer's GitHub repository for the package;
forthcoming package updates would remain in a ``development'' branch in
the same GitHub repository. Although less relevant to new R users, this
means that you can access and test upcoming versions of a package
through GitHub, which may be particularly helpful if you're developing a
package with dependencies and want to know if upcoming changes to a
package that your code depends on will cause issues. Lastly, there are
some instances where a packages is posted on CRAN---not only because
CRAN is its main repository, but also to ease its connection with base
R---then also published in a specialized repository like rOpenSci. As
part of their review process, rOpenSci asks authors of submitted
packages if they plan to publish on CRAN (or Bioconductor) and many
times packages become available in both places. If you want a package
that more or less guarantees that it does what it says it will do,
consider packages that have been peer reviewed. Knowing this, you might
start by searching ROpenSci's repository website rather than CRAN Task
Views or biocViews.

\hypertarget{rule-4-explore-the-availability-and-quality-of-help}{%
\section{Rule 4: Explore the availability and quality of
help}\label{rule-4-explore-the-availability-and-quality-of-help}}

There has been a call for the development of centralized resources in
statistical computing to enable a common understanding of software
quality and reliability: software information specified in publications,
domain-specific semantic dictionaries, and a single metadata resource
for statistical software {[}10{]}. No such resources have been
consolidated to serve these purposes and, given the decentralized nature
of today's information society, it is questionable whether they will
emerge. Current sources of information related to R packages are
dispersed and plentiful. On one hand, this allows users to explore
diverse solutions and discover new tools; on the other, not knowing
where to find help can lead to inefficient and ineffectual roundabouts.
Clearly, not all package resources share the same level of quality and
the fact that there are many resources in aggregate does not imply that
every package is associated with the same availability of resources.
While all R packages warrant some minimal standard of documentation,
beginners and users of complex packages might desire more.

Help for R packages can be broadly split into two categories: help
documentation that comes embedded in the source code of the package
(i.e., ``ships'' with the package) and documentation or other resources
that are available more broadly, beyond the package's source code. The
documentation that ships with the package can include help files,
vignettes, README files, {[}others? does ``pkgdown'' content get
included in the package source code? how about other tutorials? May be
worth checking the ``R Packages'' book on this{]}. Broader help
resources can include webpages or online books whose content is hosted
outside the package source, online tutorials, mailing lists, and large
and vocal communities of users.

At the most granular level, each R package should come with a helpfile
for each function that the package maintainer intends to be directly
used (as compared to smaller help functions that might be included
solely to be used within other package functions). This rule is enforced
for packages on some repositories, like CRAN, but might not be followed
for packages posted on less restrictive platforms, like GitHub. For each
function in the package, you can access its helpfile from the console by
typing \texttt{?} followed by the function's name (e.g., to get the
helpfile for the \texttt{mean} function, type \texttt{?mean}). Function
helpfiles include information about how to use the function, including
the parameters that can be set, the format that should be used in
setting arguments for these parameters, a description of the file, and
even references to papers or other resources on which the algorithm
underlying the function was based. They often contain executable
examples that you can use to get a feel for how to use the function. If
you would like to check all help documentation for a package, you can
access an index of help pages from the console via
\texttt{help(package\ =\ "...")}. The extent and quality of content in
these helpfiles will vary; ideally, packages should have thorough
documentation at the function level. If the help pages alone leave you
wanting, the package likely does not have further (quality)
documentation and therefore should not be your first choice, if
comparable options exist. In short, if the developer cannot initially
communicate how their tool works, then you may not want to use it in
your kitchen (read: if the instruction manual is useless, do not use the
blender).

Fortunately, plenty of R packages include additional documentation
beyond this minimum. The documentation that accompanies functions within
packages is critical; the fact that anyone can read the documentation
anytime and use it to guide their own work facilitates extensibility.
RDocumentation is a searchable
\href{https://www.rdocumentation.org/}{website}, package
(\texttt{install.packages("RDocumentation")}), and
\href{https://www.rdocumentation.org/docs/}{JSON API} for obtaining
integrated documentation for packages that are on CRAN, Bioconductor,
and GitHub. This is a subsidiary reason why packages shared on these
platforms tend to be superior (see Rule 3). RDocumentation may include:
an overview, installation instructions, examples of usage, functions,
guides, and vignettes. Most software documentation is rather technical
and extraneous to new users whereas a vignette is a practical type of
documentation in the form of a tutorial. A vignette is a detailed,
long-form document that describes the problems an R package can solve,
then illustrates applications through clear examples of code with
coordination of functions and explanations of outcomes. Packages can
have multiple vignettes; you can view or edit a specific vignette or
obtain a list of all vignettes for a package of interest via the
\texttt{vignette()} function.

Some packages are branded quite well and include a comprehensive set of
resources. Implicitly, this indicates that the authors are at least
serious about their package development, which may lead you to infer
that they know what they (and their package) are doing. Exemplary
documentation can signify an exceptional package. For instance, some
packages have websites and/or books. One popular method that developers
use to publish books about their package is through \texttt{bookdown}, a
relatively new extension of R Markdown that is structured in such a way
that integrates code, text, links, graphics, videos, and other content
in a format that can be published as a free, open, interactive, and
downloadable online book {[}40{]}. The \texttt{bookdown} package itself
has an \href{https://bookdown.org/yihui/bookdown/}{online book} that
details usage of the package {[}41{]}. \texttt{Rccp} is another package
with notable documentation and first-rate help resources; the developers
maintain both a \href{http://www.rcpp.org/}{main} and
\href{http://dirk.eddelbuettel.com/code/rcpp.html}{additional} website
with a wealth of organized information about the package and resources,
including examples, associations, publications, articles, blogs, code,
books, talks, a mailing list, and links to other resources with
\texttt{Rccp} tags. Aside from the documentation and resources from the
developer, further information about some R packages is available in
video tutorials, webinars, and code demonstrations (i.e., ``demos''). As
of RStudio v1.3, you can access tutorials powered by the \texttt{learnr}
package from the Tutorial pane in the IDE {[}42{]}. Finally, keep in
mind that RStudio creates cheatsheets of concise usage information for
popular packages through code and graphics organized by purpose.
Cheatsheets can be accessed directly via the RStudio Menu (Help
\textgreater{} Cheatsheets) or from the
\href{https://rstudio.com/resources/cheatsheets/}{RStudio website} on
which you can subscribe to cheatsheet updates and find translated
versions.

While using a package, anticipate complications beyond the scope of
documentation. In this case, you will use resources that involve
\emph{asking} for help---should the occasion arise, you want to be
assured that you will find a satisfactory answer. In the past, the
antiquated R-help mailing list was the only way to seek assistance;
since, the R community has formed, with inclusion and creative
problem-solving as hallmarks of its online presence {[}43{]}. The modern
R-help mailing list to which you can subscribe and send questions is
moderated by the R Core Development Team and includes additional facets
for major announcements about the development of R and availability of
new code (R-announce) and new or enhanced contributed packages
(R-packages) {[}44{]}. Certain packages have independent listservs;
\texttt{statnet} is an example of a suite of packages that has its own
\href{http://statnet.org/}{community listserv}. If a package has a
development repository on GitHub, check the Issues to verify that the
maintainer is responsive to posts and fixes bugs in a timely manner. In
addition, you can search discussion forums such as Stack Overflow, Cross
Validated, and Talk Stats to assess the activity associated with the
package in question. Analyses of the popularity of comparable data
analysis software in email and discussion traffic suggest that R is
rapidly becoming more prevalent and is the leading language by these
metrics {[}6,45{]}. When you encounter a problem, it is good practice to
first update the package to see if the problem is due to a bug in a
previous version---if the problem persists, seek help by finding or
posting a reproducible example {[}46{]}. Overall, avoid using a package
if the quality and quantity of related resources is lacking.

\hypertarget{rule-5-quantify-how-established-it-is}{%
\section{Rule 5: Quantify how established it
is}\label{rule-5-quantify-how-established-it-is}}

Consulting data to inform comparisons is never a bad idea; numerical
data associated with R packages will give you an impression of how
regarded the tool is and whether it has stood the test of time. Since
there are tens of thousands of R packages, you may be wondering how they
stack up in terms of popularity. On GitHub, a large number of Stars,
Forks, and Watchers associated with a package implies a substantial
following and widespread usage {[}47{]}. Likewise, the number of Google
Scholar citations is a metric of a package's impact on scientific
research and utility in research contexts (see Rule 2). RDocumentation
(see Rule 4) is rich with stats on R packages. RDocumentation hosts a
live \href{https://www.rdocumentation.org/trends}{Leaderboard} with
trends including the number of indexed packages and indexed functions,
most downloaded packages, most active maintainers, newest packages, and
newest updates. What's more, each package is assigned a percentile
rank---featured on its RDocumentation page---that quantifies the number
of times a package has been downloaded in a given month. A ranking
algorithm computes the direct, user-requested monthly downloads by
accounting for reverse dependencies (indirect downloads) so packages
that are commonly depended upon, and hence frequently downloaded, do not
skew the calculation {[}48{]}. You can research stats on corresponding
dependencies for a more holistic picture. To further determine if a
package is well-established in the R community, refer to the number of
versions and updates (more is better) as well as the date of the most
recent versions and updates (newer is better).

\hypertarget{rule-6-seek-evidence-of-peer-acceptance-and-review}{%
\section{Rule 6: Seek evidence of peer acceptance and
review}\label{rule-6-seek-evidence-of-peer-acceptance-and-review}}

Peer review is an important aspect of scientific research, not least
because it establishes scholarly credibility. You can research
information about an R package in different forms of literature and
determine the extent to which it has been validated by the scientific
community. Some journals publish articles about R packages themselves
while others feature work that used a particular package. Literature in
your field may either introduce R packages developed to solve a unique
data science problem or mention packages used during the research
process. The former may be published in the \emph{Journal of Statistical
Software}, \emph{The R Journal}, or \emph{BMC Bioinformatics}, for
example, and search queries that include \texttt{"R\ package"} with
domain keywords will narrow results. The latter requires identifying
authors who used R in their analyses; useful packages may be mentioned
in the Methods and/or References sections. You can search for packages
directly by name in Google Scholar: the \texttt{Cited\ by} link displays
the number of times a package has been cited and connects to a page with
those publications.

These packages are technically sound and have made a substantial
contribution to their fields and/or a common data science problem. In
response to the rising number of researchers creating tools and software
to work with their data, GitHub has granted developers the ability to
obtain a Digital Object Identifier (DOI) for any GitHub repository
archive so that code can be cited in academic literature {[}49{]}. If a
package has such a DOI, you can explore the network of research
associated with that package. Many R packages are associated with
content in books and series from scientific publishers such as Springer.
More directly, rOpenSci, is a unique example of an ecosystem of open
source tools with peer reviewed R packages (see Rule 3) {[}30{]}.

\hypertarget{rule-7-find-out-who-developed-it}{%
\section{Rule 7: Find out who developed
it}\label{rule-7-find-out-who-developed-it}}

Just as research is a library of shared insight, open source software is
a collection of shared tools. We care about who writes the articles we
read; we should also care about who creates the tools we use. Although R
is grounded in statistical computing and graphics, there is variation R
users' backgrounds and skills, and the same is true for R developers.
That said, the R community prides itself on embracing newcomers at all
levels of involvement. This Rule does not imply that worthwhile packages
are exclusively written by well-known authors. Rather, associations and
reputation can be a proxy for quality; in this way, the process of
evaluating and comparing R packages is no different than other
decisions. In fact, as you become more immersed in the R community, you
will find that name recognition is a factor, among many, that helps you
establish trust in certain tools more quickly than others {[}47{]}.

Information about the package authors can be found in the metadata of
the package. If you are looking at the source code of the package, there
should always be a \texttt{DESCRIPTION} file with metadata. The
\texttt{DESCRIPTION} is a succinct record of the package's purpose,
dependencies, version, date, license, associations, authors, and other
technical details. In the ``Authors'' field of this file, you determine
both the ``Maintainer'' of the package (a similar role to that of the
corresponding author for an article) as well as other authors or
contributors. For packages hosted solely on a platform like GitHub, this
might be the only way to discover a package's author; however, for
packages hosted on repositories, this package data is often also posted
somewhere online by the repository. Packages on CRAN have a page of
metadata that is accessible through a standard web address. For example,
to see the metadata page for the \texttt{ggplot2} package, you can go to
https://CRAN.R-project.org/package=ggplot2; replace the package name at
the end of the web address with the name of any other CRAN package to
get to analogous package-specific metadata. Similarly, Bioconductor
packages each have their own webpage, with some metadata at the top of
the page, including the maintainer and other authors. For example, the
page for the package \texttt{DESeq2} can be found at
https://bioconductor.org/packages/release/bioc/html/DESeq2.html.

You can assess the credibility and commitment of R package developers
through direct and indirect signals. Who made the package? Consider
whether expertise in a certain domain is vital to the design and
creation of the tool. Research the authors' associations in academia,
industry, and/or laboratories and gauge the extent to which they have a
primary role in R development. Further, you can learn more about their
experience, active contributions to the R community, and history related
to package development by exploring their profiles on GitHub, Google
Scholar, Research Gate, Twitter, or personal or package websites. If an
author has such a history, peruse their portfolio of packages to see if
any are highly regarded or recognizable. Frequent commits and effective
resolutions of GitHub Issues can reveal the authors' priorities and
commitment. If the package was developed by multiple authors, research
each of them to evaluate the robustness of the team. By extension, these
indicators of developer involvement and reputation will help you discern
whether a package is worthy of your trust or requires evaluation based
on other Rules.

\hypertarget{rule-8-see-how-its-developed}{%
\section{Rule 8: See how it's
developed}\label{rule-8-see-how-its-developed}}

You do not need to be a software engineer to identify strong package
development. Scientific software developers sometimes neglect best
practices; indeed, these shortcomings are evident in the tools they
create {[}50{]}. There are concrete ways to measure a tool's robustness
beyond whether it works for those who did not create it. R packages
often depend on other R packages; you should check the reputations of
such \emph{dependencies} when selecting a package---quality packages
will rely on a solid web of quality packages. What's more, like other
types of software, well-maintained R packages have multiple versions
corresponding to iterative releases to indicate that the package is
compatible with dependencies and loyally updated (e.g., bug fixes,
general improvements, new functionality) {[}1,24{]}. You can explore the
version history of a package to see if it is up-to-date. As a user,
there are two additional development protocols that you can further
investigate to assess the underlying stability and utility of a package:
unit tests and version control.

A responsible developer with a consistent and reproducible workflow will
implement formal testing on their code to examine expected behavior via
an automated process called unit testing {[}24,51{]}. Although
inconvenient at the outset, the developer---and by extension, the
package user---will benefit from unit testing, which results in fewer
bugs, a well-designed code structure, an efficient workflow, and robust
code that is not sensitive to major changes in the future {[}24{]}. To
alleviate the burdens of unit testing, \texttt{testthat} is a popular,
integrative R package that helps developers create reliable functions,
minimize error, and visualize progress through automatic code testing
{[}52{]}. Developers are also interested in quantifying the amount of
code in their package that has been tested. Test coverage, a measurement
of the proportion of code that has undergone unit testing, is an
objective metric for package developers, contributors, and users to
evaluate code quality. Many developers use the \texttt{covr} package to
generate reports and determine the magnitude of coverage on the
function, script, and package levels {[}53{]}. Relatedly, developers who
host their packages on GitHub, post status badges in the overview
(\texttt{README}) section of the repository webpage. GitHub badges are a
common self-imposed method to signal use of best practices and motivate
developers to produce a product that is high in quality and transparency
{[}54{]}. You may see, for example, license, dependency, or style
badges, all of which are good indicators of package caliber; however,
particular to this Rule, you should look for code coverage
(\texttt{codecov}) badges which reveal the percentage of test coverage.

As we mentioned in Rule 3, version control has an essential role in
package development and computational literacy more broadly {[}24,55{]}.
Version control is like a time capsule for your workflow because it
monitors and tracks changes to files as a project evolves, and stores
them as previous versions to be recovered if necessary. In other words,
``version control is as fundamental to programming as accurate notes
about lab procedures are to experimental science'' {[}{[}55{]}; p6{]}.
Git is a decentralized open source version control system that is useful
regardless of whether a project is independent or collaborative
{[}56{]}. GitHub works in conjunction with Git to provide a powerful
structured system to organize and manage components of a project for
others and your future self. A growing number of scientists have
research programs based in GitHub, which has become a revolutionary tool
for productive team science and distributed development efforts
{[}1,57{]}. As you may expect, Git coupled with GitHub is the version
control duo of choice among serious R package developers {[}24{]}. Thus,
if the package you are interested in using is among the thousands hosted
on GitHub, this is evidence that the developer is at least committed to
a logical, open, and reproducible workflow, suggestive of more time
spent designing their tool.

\hypertarget{rule-9-put-it-to-the-test}{%
\section{Rule 9: Put it to the test}\label{rule-9-put-it-to-the-test}}

If you are unable to decide whether to use a package based on prior
Rules, test it out. Similarly, if you have narrowed your options, work
with each to highlight differences. Exploring the package and engaging
in trial and error using your skills in context of your goal will
illuminate technical details and solidify any doubts. Note, in the case
that the package you want to try has been shared as a zipped file, you
can use a GitHub mirror of CRAN via
\texttt{devtools::install\_github("username/reponame")} as an
alternative to downloading a large or potentially corrupted zipped file.

At this point, what you have learned about the package should be quite
helpful. If the development and documentation are sound, the package
should come with a test script or working example that you can run after
installation {[}50{]}. Vignettes include many common data science
problems with solutions; you can run the code examples, tweak them, and
compare the outputs. In general, it is essential to know the behavior of
different functions within a package, how they interact, and how outputs
respond to changes in inputs. Suppose you are testing a package with
sparse documentation such that function descriptions often include
``\texttt{\ldots{}}'' and the argument descriptions seem incomplete.
This will be problematic if making a reasonable change to an argument
results in an incomprehensible error for which you cannot find help.
When this happens, you may not want to use the package for your task.

Sometimes packages do not interact well with other packages; a recipe
prepared with an odd combination of tools will not turn out. If you are
interested in working with a certain package but are already using other
packages in your workflow, you will need to verify that they work
together. More precisely, you should check the \emph{interoperability}
of all the packages you want to use. A given package may be highly
specialized and incompatible with certain packages in general, or simply
have a few tolerable quirks for which you can develop workarounds. There
are some packages that are masterful at doing what they are made to do,
yet incongruous with other packages. Such packages might, for example,
use S3 or S4 objects, which are two main approaches developers use to
implement object-oriented programming in R. Many packages for spatial
analysis as well as those from Bioconductor tend to use S4 objects to
represent data {[}7{]}. On the other hand, the \texttt{tidyverse}, a
unified suite of packages with an grammatical structure employed within
a ``pipeline'', expects data frame objects {[}58{]}. Thus, when you are
working in the \texttt{tidyverse}, you cannot incorporate S3 and S4
objects into the framework unless their corresponding functions are the
final step in the pipeline. The \texttt{broom} package, and the
bioinformatics analog, \texttt{biobroom}, aim to alleviate these
disruptions by converting untidy objects into tidy data, thereby making
it easier to integrate statistical functions into the structure of the
\texttt{tidyverse} workflow {[}59,60{]}. Furthermore, the \texttt{caret}
package facilitates interoperability for machine learning packages by
providing a uniform interface for modeling with various algorithms from
different packages that would otherwise have independent syntax
{[}61{]}.

\hypertarget{rule-10-develop-your-own-package}{%
\section{Rule 10: Develop your own
package}\label{rule-10-develop-your-own-package}}

Alternative solutions can be sought when a package to solve your data
science problem is nonexistent. An R package is the fundamental unit of
shareable code; rather than exclusively being a user of packages, you
can create them---more easily than you may think {[}24{]}. The reasons
why you might want to create a package are abundant, including
necessity, innovation, standardization, automation, specialty,
containment, organization, sharing, collaboration, and extensibility.
The essence of an R package is a self-contained piece of statistical
knowledge that can be used in combination with other self-contained
pieces of statistical knowledge of different shapes and sizes; the
uniquely structured functions within a package help us implement that
knowledge and weave it into novel scientific work.

Whatever your motivation, packages are simply tools; you can create a
package out of any collection of specialty functions. Packages need not
be formal nor entirely cohesive. For instance, personal R packages such
as \texttt{Hmisc} and \texttt{broman}, are comprised of miscellaneous
functions which the creator has developed and frequently uses
{[}62,63{]}. Functions are necessary for efficiency and warranted when
you repetitiously copy and paste your code while making slight
modifications after each iteration {[}46{]}. The concept of personal R
packages demonstrates a unique purpose for packages beyond the
conventional. R packages are not solely reserved for specific tasks with
comprehensive methods; rather, package development can help you learn
how to apply proper coding techniques to writing functions and
documentation with reproducibility and collaboration in mind {[}64{]}.

Although you may not anticipate that anyone else will use your tools,
following best practices for package development will yield more
favorable outcomes. As a consumer of shared packages, you know the
inherent benefits of robust software development relative to the quality
of code, data, documentation, versions, and tests {[}50{]}. Similarly,
creating a valuable package for personal use requires consideration for
your future self and anticipation of distributing your code, should the
need arise. Use version control and take advantage of existing
resources. Indeed, there are R packages that aid in package development
(e.g., \texttt{devtools}, \texttt{usethis}, \texttt{testthat},
\texttt{roxygen2}, \texttt{rlang}, \texttt{drat}) {[}32,38,65--68{]}. In
the case of collaboration, the R project within RStudio IDE, is
compatible with distributed development---a feature that couples well
with version control. There is no lack of effective organizational
frameworks to reference in the open source R community; in fact,
repositories for many exemplary packages are available on GitHub. We
recommend consulting resources authored by expert R developers including
\href{https://r-pkgs.org/}{\emph{R Packages}} by Hadley Wickham and the
official manual,
\href{https://cran.r-project.org/doc/manuals/r-release/R-exts.html}{\emph{Writing
R Extensions}}, from CRAN {[}24,69{]}.

\hypertarget{conclusion}{%
\section{Conclusion}\label{conclusion}}

R packages are a defining feature of the language insofar as many are
robust, user-friendly, and richly extend R's core functionality. Some of
the most prominent R packages are a result of the developer abstracting
common elements of a data science problem into a workflow that can be
shared and accompanied by thorough descriptions of the process and
purpose. In this way, R packages have effectively transformed how we
interact with data in the modern day in, perhaps, a more impactful
manner than several revered contributions to theoretical statistics
{[}4{]}. Packages greatly enhance the user experience and enable you to
be more efficient and effective at learning from data, regardless of
prior experience.

Nonetheless, the sheer quantity and potential complexity of available R
packages can undermine their collective benefits. Finding and choosing
packages, particularly for beginners, can be daunting and difficult. R
users often struggle to sift through the tools at their disposal and
wonder how to distinguish appropriate usage. These ten simple rules for
navigating the shared code in the R community are intended to serve as a
valuable page in your computing cookbook---one that will evolve into
intuition and yet remain a reliable reference. May searching for and
selecting proper tools no longer spoil your appetite and dissuade you
from discovering, trying, creating, and sharing new recipes.

\hypertarget{table-1-general-packages}{%
\section{Table 1 (general packages)}\label{table-1-general-packages}}

\begin{Shaded}
\begin{Highlighting}[]
\KeywordTok{library}\NormalTok{(kableExtra)}
\KeywordTok{library}\NormalTok{(knitr)}
\end{Highlighting}
\end{Shaded}

\begin{Shaded}
\begin{Highlighting}[]
\CommentTok{# general packages data}
\NormalTok{gen_pkgs <-}\StringTok{ }\KeywordTok{data.frame}\NormalTok{(}
  \DataTypeTok{Package =} \KeywordTok{c}\NormalTok{(}\StringTok{"readr"}\NormalTok{, }
              \StringTok{"dplyr"}\NormalTok{,}
              \StringTok{"tidyr"}\NormalTok{,}
              
              \StringTok{"broom[note]"}\NormalTok{,}
              \StringTok{"purrr"}\NormalTok{,}
              \StringTok{"caret"}\NormalTok{, }
              
              \StringTok{"ggplot2"}\NormalTok{, }
              \StringTok{"kableExtra"}\NormalTok{, }
              \StringTok{"rmarkdown"}\NormalTok{),}
  
  \DataTypeTok{Description =} \KeywordTok{c}\NormalTok{(}\StringTok{"Read rectangular data (e.g., csv, tsv, and fwf)"}\NormalTok{, }
                  \StringTok{"Grammar of tidy data manipulation"}\NormalTok{,}
                  \StringTok{"Tidy messy data"}\NormalTok{,}
                  
                  \StringTok{"Tidy model output; convert statistical objects into tidy tibbles"}\NormalTok{, }
                  \StringTok{"Functional programming tools for functions and vectors"}\NormalTok{,}
                  \StringTok{"Framework for predictive modeling"}\NormalTok{,}
                  
                  \StringTok{"System for data visualization"}\NormalTok{, }
                  \StringTok{"Build complex tables and manipulate styles"}\NormalTok{,}
                  \StringTok{"Authoring framework for data science and reproducible research"}\NormalTok{),}
  
  \DataTypeTok{Year =} \KeywordTok{c}\NormalTok{(}\StringTok{"2015"}\NormalTok{,}
           \StringTok{"2014"}\NormalTok{,}
           \StringTok{"2014"}\NormalTok{,}
  
           \StringTok{"2014"}\NormalTok{,}
           \StringTok{"2015"}\NormalTok{,}
           \StringTok{"2007"}\NormalTok{,}
                    
           \StringTok{"2007"}\NormalTok{,}
           \StringTok{"2017"}\NormalTok{, }
           \StringTok{"2014"}\NormalTok{),}
  
  \DataTypeTok{Author =} \KeywordTok{c}\NormalTok{(}\StringTok{"Wickham et al."}\NormalTok{,}
             \StringTok{"Wickham et al."}\NormalTok{, }
             \StringTok{"Wickham & Henry"}\NormalTok{,}
  
             \StringTok{"Robinson & Hayes"}\NormalTok{,}
             \StringTok{"Henry & Wickham"}\NormalTok{,}
             \StringTok{"Kuhn"}\NormalTok{, }
                    
             \StringTok{"Wickham et al."}\NormalTok{,}
             \StringTok{"Zhu"}\NormalTok{, }
             \StringTok{"Allaire et al."}\NormalTok{),}
  
  \DataTypeTok{Documentation =} \KeywordTok{c}\NormalTok{(}\StringTok{"https://readr.tidyverse.org/"}\NormalTok{, }
                    \StringTok{"https://dplyr.tidyverse.org/"}\NormalTok{, }
                    \StringTok{"https://tidyr.tidyverse.org/"}\NormalTok{,}
  
                    \StringTok{"https://broom.tidymodels.org/"}\NormalTok{,}
                    \StringTok{"https://purrr.tidyverse.org/"}\NormalTok{,}
                    \StringTok{"https://topepo.github.io/caret/index.html"}\NormalTok{, }
                    
                    \StringTok{"https://ggplot2.tidyverse.org/"}\NormalTok{,}
                    \StringTok{"https://haozhu233.github.io/kableExtra/"}\NormalTok{, }
                    \StringTok{"https://rmarkdown.rstudio.com/lesson-1.html"}\NormalTok{)}
\NormalTok{)}
\end{Highlighting}
\end{Shaded}

\begin{Shaded}
\begin{Highlighting}[]
\CommentTok{# general packages table}
\KeywordTok{kable}\NormalTok{(gen_pkgs, }\DataTypeTok{format =} \StringTok{"latex"}\NormalTok{, }\DataTypeTok{booktabs =} \OtherTok{TRUE}\NormalTok{) }\OperatorTok{%>%}
\StringTok{  }\CommentTok{# scale}
\StringTok{  }\KeywordTok{kable_styling}\NormalTok{(}\DataTypeTok{latex_options =} \StringTok{"scale_down"}\NormalTok{) }\OperatorTok{%>%}
\StringTok{  }\CommentTok{# separate rows by category}
\StringTok{  }\KeywordTok{pack_rows}\NormalTok{(}\StringTok{"Data Manipulation"}\NormalTok{, }\DecValTok{1}\NormalTok{, }\DecValTok{3}\NormalTok{) }\OperatorTok{%>%}\StringTok{ }
\StringTok{  }\KeywordTok{pack_rows}\NormalTok{(}\StringTok{"Statistical Modeling"}\NormalTok{, }\DecValTok{4}\NormalTok{, }\DecValTok{6}\NormalTok{) }\OperatorTok{%>%}\StringTok{ }
\StringTok{  }\KeywordTok{pack_rows}\NormalTok{(}\StringTok{"Data Visualization"}\NormalTok{, }\DecValTok{7}\NormalTok{, }\DecValTok{9}\NormalTok{) }\OperatorTok{%>%}
\StringTok{  }\CommentTok{# column wrap}
\StringTok{  }\KeywordTok{column_spec}\NormalTok{(}\DecValTok{1}\NormalTok{, }\DataTypeTok{width =} \StringTok{"10em"}\NormalTok{) }\OperatorTok{%>%}\StringTok{ }
\StringTok{  }\KeywordTok{column_spec}\NormalTok{(}\DecValTok{2}\NormalTok{, }\DataTypeTok{width =} \StringTok{"20em"}\NormalTok{) }\OperatorTok{%>%}\StringTok{ }
\StringTok{  }\CommentTok{# bold column names}
\StringTok{  }\KeywordTok{row_spec}\NormalTok{(}\DecValTok{0}\NormalTok{, }\DataTypeTok{bold =}\NormalTok{ T) }\OperatorTok{%>%}\StringTok{ }
\StringTok{  }\KeywordTok{add_footnote}\NormalTok{(}\StringTok{"See the biobroom analog in Bioconductor"}\NormalTok{,}
               \DataTypeTok{notation =} \StringTok{"symbol"}\NormalTok{)}
\end{Highlighting}
\end{Shaded}

\begin{table}[H]
\centering
\resizebox{\linewidth}{!}{
\begin{tabular}{>{\raggedright\arraybackslash}p{10em}>{\raggedright\arraybackslash}p{20em}lll}
\toprule
\textbf{Package} & \textbf{Description} & \textbf{Year} & \textbf{Author} & \textbf{Documentation}\\
\midrule
\addlinespace[0.3em]
\multicolumn{5}{l}{\textbf{Data Manipulation}}\\
\hspace{1em}readr & Read rectangular data (e.g., csv, tsv, and fwf) & 2015 & Wickham et al. & https://readr.tidyverse.org/\\
\hspace{1em}dplyr & Grammar of tidy data manipulation & 2014 & Wickham et al. & https://dplyr.tidyverse.org/\\
\hspace{1em}tidyr & Tidy messy data & 2014 & Wickham \& Henry & https://tidyr.tidyverse.org/\\
\addlinespace[0.3em]
\multicolumn{5}{l}{\textbf{Statistical Modeling}}\\
\hspace{1em}broom\textsuperscript{*} & Tidy model output; convert statistical objects into tidy tibbles & 2014 & Robinson \& Hayes & https://broom.tidymodels.org/\\
\hspace{1em}purrr & Functional programming tools for functions and vectors & 2015 & Henry \& Wickham & https://purrr.tidyverse.org/\\
\hspace{1em}caret & Framework for predictive modeling & 2007 & Kuhn & https://topepo.github.io/caret/index.html\\
\addlinespace[0.3em]
\multicolumn{5}{l}{\textbf{Data Visualization}}\\
\hspace{1em}ggplot2 & System for data visualization & 2007 & Wickham et al. & https://ggplot2.tidyverse.org/\\
\hspace{1em}kableExtra & Build complex tables and manipulate styles & 2017 & Zhu & https://haozhu233.github.io/kableExtra/\\
\hspace{1em}rmarkdown & Authoring framework for data science and reproducible research & 2014 & Allaire et al. & https://rmarkdown.rstudio.com/lesson-1.html\\
\bottomrule
\multicolumn{5}{l}{\textsuperscript{*} See the biobroom analog in Bioconductor}\\
\end{tabular}}
\end{table}

\begin{Shaded}
\begin{Highlighting}[]
\CommentTok{## trying to separate color; striped by group}
\CommentTok{#  row_spec(1:3 - 1, extra_latex_after = "\textbackslash{}\textbackslash{}rowcolor\{gray!6\}")}
\CommentTok{#  row_spec(0:3, extra_latex_after = "\textbackslash{}\textbackslash{}rowcolor\{orange!6\}") %>% }
\CommentTok{#  row_spec(4:6, extra_latex_after = "\textbackslash{}\textbackslash{}rowcolor\{gray!6\}") %>% }
\CommentTok{#  row_spec(7:11, extra_latex_after = "\textbackslash{}\textbackslash{}rowcolor\{gray!6\}")}

\CommentTok{## QUESTIONS }
\CommentTok{# Code font for package names in " "? \textbackslash{}texttt\{\}?}
\CommentTok{# How do you repeat same symbol on multiple items with one footnote?}
\CommentTok{# How do you separate colors and stripe by group?}
\CommentTok{# Add title}
\CommentTok{# Add caption}
\CommentTok{# Cite packages in bib and add references in table?}
\CommentTok{# Embed url link to package documentation? Do we want to link cheatsheets?}
\CommentTok{# How do you add link/reference to Table 1 in text in the template?}
\CommentTok{# How do you hide code for table in knitted pdf...include=FALSE errors?}
\CommentTok{# Title for column 2: description/purpose/usage?}
\CommentTok{# Length of description/purpose/usage for each package?}
\end{Highlighting}
\end{Shaded}

\hypertarget{supporting-information}{%
\section{Supporting information}\label{supporting-information}}

Do we need to include any supporting information?

\hypertarget{acknowledgements}{%
\section{Acknowledgements}\label{acknowledgements}}

{[}Acknowledgment of people who have helped; suggestions from
colleagues, etc.{]}

\hypertarget{references}{%
\section*{References}\label{references}}
\addcontentsline{toc}{section}{References}

\hypertarget{refs}{}
\leavevmode\hypertarget{ref-perez2016}{}%
1. Perez-Riverol Y, Gatto L, Wang R, Sachsenberg T, Uszkoreit J, Veiga
Leprevost F da, et al. Ten simple rules for taking advantage of git and
github. PLoS computational biology. Public Library of Science; 2016;12.

\leavevmode\hypertarget{ref-peng2011}{}%
2. Peng RD. Reproducible research in computational science. Science.
American Association for the Advancement of Science; 2011;334:
1226--1227.

\leavevmode\hypertarget{ref-goodman2016}{}%
3. Goodman SN, Fanelli D, Ioannidis JP. What does research
reproducibility mean? Science translational medicine. American
Association for the Advancement of Science; 2016;8: 341ps12--341ps12.

\leavevmode\hypertarget{ref-donoho2017}{}%
4. Donoho D. 50 years of data science. Journal of Computational and
Graphical Statistics. Taylor \& Francis; 2017;26: 745--766.

\leavevmode\hypertarget{ref-stodden2013}{}%
5. Stodden V, Miguez S. Best practices for computational science:
Software infrastructure and environments for reproducible and extensible
research. Available at SSRN 2322276. 2013;

\leavevmode\hypertarget{ref-robinson2017}{}%
6. Robinson D. The impressive growth of r {[}Internet{]}. Stack
Overflow; 2017. Available:
\url{https://stackoverflow.blog/2017/10/10/impressive-growth-r/}

\leavevmode\hypertarget{ref-gentleman2004}{}%
7. Gentleman RC, Carey VJ, Bates DM, Bolstad B, Dettling M, Dudoit S, et
al. Bioconductor: Open software development for computational biology
and bioinformatics. Genome biology. Springer; 2004;5: R80.

\leavevmode\hypertarget{ref-holmes2018}{}%
8. Holmes S, Huber W. Modern statistics for modern biology
{[}Internet{]}. Cambridge University Press; 2018. Available:
\url{https://web.stanford.edu/class/bios221/book/index.html}

\leavevmode\hypertarget{ref-Rproject2020}{}%
9. Team RC. The r project for statistical computing {[}Internet{]}. The
R Foundation; 2020. Available: \url{https://www.r-project.org/}

\leavevmode\hypertarget{ref-hornik2012}{}%
10. Hornik K. Are there too many r packages? Austrian Journal of
Statistics. 2012;41: 59--66.

\leavevmode\hypertarget{ref-dplyr}{}%
11. Wickham H, François R, Henry L, Müller K. Dplyr: A grammar of data
manipulation {[}Internet{]}. 2020. Available:
\url{https://CRAN.R-project.org/package=dplyr}

\leavevmode\hypertarget{ref-tidyr}{}%
12. Wickham H, Henry L. Tidyr: Tidy messy data {[}Internet{]}. 2020.
Available: \url{https://CRAN.R-project.org/package=tidyr}

\leavevmode\hypertarget{ref-twitteR}{}%
13. Gentry J. TwitteR: R based twitter client {[}Internet{]}. 2015.
Available: \url{https://CRAN.R-project.org/package=twitteR}

\leavevmode\hypertarget{ref-mailR}{}%
14. Premraj R. MailR: A utility to send emails from r {[}Internet{]}.
2015. Available: \url{https://CRAN.R-project.org/package=mailR}

\leavevmode\hypertarget{ref-phonenumber}{}%
15. Myles S. Phonenumber: Convert letters to numbers and back as on a
telephone keypad {[}Internet{]}. 2015. Available:
\url{https://CRAN.R-project.org/package=phonenumber}

\leavevmode\hypertarget{ref-bioCproject}{}%
16. Huber W, Carey VJ, Gentleman R, Anders S, Carlson M, Carvalho BS, et
al. Orchestrating high-throughput genomic analysis with Bioconductor.
Nature Methods. 2015;12: 115--121. Available:
\url{http://www.nature.com/nmeth/journal/v12/n2/full/nmeth.3252.html}

\leavevmode\hypertarget{ref-zeileis2005}{}%
17. Zeileis A. CRAN task views. R News. 2005;5: 39--40.

\leavevmode\hypertarget{ref-smith2017}{}%
18. Smith D. The r community is one of r's best features {[}Internet{]}.
Revolutions. Microsoft; 2017. Available:
\url{https://blog.revolutionanalytics.com/2017/06/r-community.html}

\leavevmode\hypertarget{ref-ellis2017}{}%
19. Ellis SE. Hey! You there! You are welcome here {[}Internet{]}.
rOpenSci. NumFOCUS; 2017. Available:
\url{https://ropensci.org/blog/2017/06/23/community/}

\leavevmode\hypertarget{ref-rickert2018}{}%
20. Rickert J. What makes a great r package? {[}Internet{]}. RStudio;
2018. Available:
\url{https://rstudio.com/resources/rstudioconf-2018/what-makes-a-great-r-package-joseph-rickert/}

\leavevmode\hypertarget{ref-randomforest}{}%
21. Liaw A, Wiener M. Classification and regression by randomForest. R
News. 2002;2: 18--22. Available:
\url{https://CRAN.R-project.org/doc/Rnews/}

\leavevmode\hypertarget{ref-rayshader}{}%
22. Morgan-Wall T. Rayshader: Create maps and visualize data in 2D and
3D {[}Internet{]}. 2020. Available:
\url{https://CRAN.R-project.org/package=rayshader}

\leavevmode\hypertarget{ref-broman2020}{}%
23. Broman K. Getting your r package on cran {[}Internet{]}. 2020.
Available: \url{https://kbroman.org/pkg_primer/pages/cran.html}

\leavevmode\hypertarget{ref-wickham2015}{}%
24. Wickham H. R packages: Organize, test, document, and share your
code. "O'Reilly Media, Inc."; 2015.

\leavevmode\hypertarget{ref-cranpolicy2020}{}%
25. CRAN repository policy {[}Internet{]}. The R Foundation; 2020.
Available:
\url{https://cran.r-project.org/web/packages/policies.html\#Submission}

\leavevmode\hypertarget{ref-weathermetrics}{}%
26. Anderson GB, Bell ML, Peng RD. Methods to calculate the heat index
as an exposure metric in environmental health research. Environemtental
Health Perspectives. 2013;121: 1111--1119. Available:
\url{http://ehp.niehs.nih.gov/1206273/}

\leavevmode\hypertarget{ref-bioconductor}{}%
27. Bioconductor: Open source software for bioinformatics
{[}Internet{]}. Bioconductor; 2020. Available:
\url{https://www.bioconductor.org/}

\leavevmode\hypertarget{ref-falcon2007}{}%
28. Falcon S, Morgan M, Gentleman R. An introduction to bioconductor's
expressionset class. 2007.

\leavevmode\hypertarget{ref-biocpkgsub2020}{}%
29. Package submission {[}Internet{]}. Bioconductor; 2020. Available:
\url{https://www.bioconductor.org/developers/package-submission/}

\leavevmode\hypertarget{ref-ropensci2020}{}%
30. Transforming science through open data and software {[}Internet{]}.
rOpenSci; 2020. Available: \url{https://ropensci.org/}

\leavevmode\hypertarget{ref-ropenscidev2020}{}%
31. rOpenSci, Anderson B, Chamberlain S, Krystalli A, Mullen L, Ram K,
et al. Ropensci/dev\_guide: Fourth release {[}Internet{]}. Zenodo; 2020.
doi:\href{https://doi.org/10.5281/zenodo.3749013}{10.5281/zenodo.3749013}

\leavevmode\hypertarget{ref-drat}{}%
32. Carl Boettiger DE with contributions by, Fultz N, Gibb S, Gillespie
C, Górecki J, Jones M, et al. Drat: 'Drat' r archive template
{[}Internet{]}. 2020. Available:
\url{https://CRAN.R-project.org/package=drat}

\leavevmode\hypertarget{ref-anderson2017}{}%
33. Anderson GB, Eddelbuettel D. Hosting data packages via drat: A case
study with hurricane exposure data. R Journal. 2017;9.

\leavevmode\hypertarget{ref-theussl2009}{}%
34. Theußl S, Zeileis A. Collaborative software development using
r-forge. Special invited paper on" the future of r". The R Journal. The
R Foundation for Statistical Computing; 2009;1: 9--14.

\leavevmode\hypertarget{ref-lang2000}{}%
35. Lang DT. The omegahat environment: New possibilities for statistical
computing. Journal of Computational and Graphical Statistics. Taylor \&
Francis; 2000;9: 423--451.

\leavevmode\hypertarget{ref-mcelreath2020}{}%
36. McElreath R. Statistical rethinking: A bayesian course with examples
in r and stan. CRC press; 2020.

\leavevmode\hypertarget{ref-decan2015}{}%
37. Decan A, Mens T, Claes M, Grosjean P. On the development and
distribution of r packages: An empirical analysis of the r ecosystem.
Proceedings of the 2015 european conference on software architecture
workshops. 2015. pp. 1--6.

\leavevmode\hypertarget{ref-devtools}{}%
38. Wickham H, Hester J, Chang W. Devtools: Tools to make developing r
packages easier {[}Internet{]}. 2020. Available:
\url{https://CRAN.R-project.org/package=devtools}

\leavevmode\hypertarget{ref-decan2016}{}%
39. Decan A, Mens T, Claes M, Grosjean P. When github meets cran: An
analysis of inter-repository package dependency problems. 2016 ieee 23rd
international conference on software analysis, evolution, and
reengineering (saner). IEEE; 2016. pp. 493--504.

\leavevmode\hypertarget{ref-bookdown}{}%
40. Xie Y. Bookdown: Authoring books and technical documents with r
markdown {[}Internet{]}. 2020. Available:
\url{https://CRAN.R-project.org/package=bookdown}

\leavevmode\hypertarget{ref-xie2016}{}%
41. Xie Y. Bookdown: Authoring books and technical documents with r
markdown. Chapman; Hall/CRC; 2016.

\leavevmode\hypertarget{ref-ushey2020}{}%
42. Ushey K. RStudio 1.3 preview: Integrated tutorials {[}Internet{]}.
RStudio; 2020. Available:
\url{https://blog.rstudio.com/2020/02/25/rstudio-1-3-integrated-tutorials/}

\leavevmode\hypertarget{ref-chase2020}{}%
43. Chase W. Dataviz and the 20th anniversary of r, an interview with
hadley wickham {[}Internet{]}. Medium; 2020. Available:
\url{https://medium.com/nightingale/dataviz-and-the-20th-anniversary-of-r-an-interview-with-hadley-wickham-ea245078fc8a}

\leavevmode\hypertarget{ref-Rmail2020}{}%
44. Team RC. Mailing lists {[}Internet{]}. The R Foundation; 2020.
Available: \url{https://www.r-project.org/mail.html}

\leavevmode\hypertarget{ref-muenchen2012}{}%
45. Muenchen RA. The popularity of data analysis software. URL
http://r4statscom/popularity. 2012;

\leavevmode\hypertarget{ref-wickham2014}{}%
46. Wickham H. Advanced r. CRC press; 2014.

\leavevmode\hypertarget{ref-leek2015}{}%
47. Leek J. How i decide when to trust an r package {[}Internet{]}.
2015. Available:
\url{https://simplystatistics.org/2015/11/06/how-i-decide-when-to-trust-an-r-package/}

\leavevmode\hypertarget{ref-vannoorenberghe2017}{}%
48. Vannoorenberghe L. RDocumentation: Scoring and ranking
{[}Internet{]}. DataCamp; 2017. Available:
\url{https://www.datacamp.com/community/blog/rdocumentation-ranking-scoring}

\leavevmode\hypertarget{ref-smith2014}{}%
49. Smith A. Improving github for science {[}Internet{]}. GitHub, Inc.
2014. Available:
\url{https://github.blog/2014-05-14-improving-github-for-science/}

\leavevmode\hypertarget{ref-taschuk2017}{}%
50. Taschuk M, Wilson G. Ten simple rules for making research software
more robust. PLoS computational biology. Public Library of Science;
2017;13.

\leavevmode\hypertarget{ref-hester2020}{}%
51. Hester J. How does covr work anyway? {[}Internet{]}. The R
Foundation; 2020. Available:
\url{https://cran.r-project.org/web/packages/covr/vignettes/how_it_works.html}

\leavevmode\hypertarget{ref-wickham2011}{}%
52. Wickham H. Testthat: Get started with testing. The R Journal.
2011;3: 5--10. Available:
\url{https://journal.r-project.org/archive/2011-1/RJournal_2011-1_Wickham.pdf}

\leavevmode\hypertarget{ref-covr}{}%
53. Hester J. Covr: Test coverage for packages {[}Internet{]}. 2020.
Available: \url{https://CRAN.R-project.org/package=covr}

\leavevmode\hypertarget{ref-barts2018}{}%
54. Barts C. How to use github badges to stop feeling like a noob
{[}Internet{]}. freeCodeCamp; 2018. Available:
\url{https://www.freecodecamp.org/news/how-to-use-badges-to-stop-feeling-like-a-noob-d4e6600d37d2/}

\leavevmode\hypertarget{ref-wilson2006}{}%
55. Wilson GV. Where's the real bottleneck in scientific computing?
American Scientist. 2006;94: 5.

\leavevmode\hypertarget{ref-bryan2018}{}%
56. Bryan J. Excuse me, do you have a moment to talk about version
control? The American Statistician. Taylor \& Francis; 2018;72: 20--27.

\leavevmode\hypertarget{ref-perkel2016}{}%
57. Perkel J. Democratic databases: Science on github. Nature News.
2016;538: 127.

\leavevmode\hypertarget{ref-tidyverse}{}%
58. Wickham H, Averick M, Bryan J, Chang W, McGowan LD, François R, et
al. Welcome to the tidyverse. Journal of Open Source Software. 2019;4:
1686.
doi:\href{https://doi.org/10.21105/joss.01686}{10.21105/joss.01686}

\leavevmode\hypertarget{ref-broom}{}%
59. Robinson D, Hayes A. Broom: Convert statistical analysis objects
into tidy tibbles {[}Internet{]}. 2020. Available:
\url{https://CRAN.R-project.org/package=broom}

\leavevmode\hypertarget{ref-biobroom}{}%
60. Andrew J. Bass SL David G. Robinson. Biobroom: Turn bioconductor
objects into tidy data frames {[}Internet{]}. 2020. Available:
\url{https://github.com/StoreyLab/biobroom}

\leavevmode\hypertarget{ref-caret}{}%
61. Kuhn M. Caret: Classification and regression training
{[}Internet{]}. 2020. Available:
\url{https://CRAN.R-project.org/package=caret}

\leavevmode\hypertarget{ref-Hmisc}{}%
62. Harrell Jr FE, Charles Dupont, others. Hmisc: Harrell miscellaneous
{[}Internet{]}. 2020. Available:
\url{https://CRAN.R-project.org/package=Hmisc}

\leavevmode\hypertarget{ref-broman}{}%
63. Broman KW. Broman: Karl broman's r code {[}Internet{]}. 2020.
Available: \url{https://CRAN.R-project.org/package=broman}

\leavevmode\hypertarget{ref-parker2013}{}%
64. Parker H. Personal r packages {[}Internet{]}. 2013. Available:
\url{https://hilaryparker.com/2013/04/03/personal-r-packages/}

\leavevmode\hypertarget{ref-usethis}{}%
65. Wickham H, Bryan J. Usethis: Automate package and project setup
{[}Internet{]}. 2019. Available:
\url{https://CRAN.R-project.org/package=usethis}

\leavevmode\hypertarget{ref-testthat}{}%
66. Wickham H. Testthat: Get started with testing. The R Journal.
2011;3: 5--10. Available:
\url{https://journal.r-project.org/archive/2011-1/RJournal_2011-1_Wickham.pdf}

\leavevmode\hypertarget{ref-roxygen2}{}%
67. Wickham H, Danenberg P, Csárdi G, Eugster M. Roxygen2: In-line
documentation for r {[}Internet{]}. 2020. Available:
\url{https://CRAN.R-project.org/package=roxygen2}

\leavevmode\hypertarget{ref-rlang}{}%
68. Henry L, Wickham H. Rlang: Functions for base types and core r and
'tidyverse' features {[}Internet{]}. 2020. Available:
\url{https://CRAN.R-project.org/package=rlang}

\leavevmode\hypertarget{ref-Rcore2020}{}%
69. Team RC. Writing r extensions {[}Internet{]}. The R Foundation;
2020. Available:
\url{https://cran.r-project.org/doc/manuals/R-exts.html}

\nolinenumbers


\end{document}

